% !TEX root = ./main.tex

\fdusetup{
	style = {font-size = -4,
			bib-backend =bibtex,
            bib-resource = {main.bib},
			bib-style = numerical
   },
	info = {
		title = {可无限分支的$N$-分支随机游走的渐进速率性质},
		title* = {Asymptotic Speed of the $N$-Branching Random Walk with Infinite Offspring},
		author = {王逸扬},
		author* = {Wang Yiyang},
		supervisor = {Élie Aïdékon 教~授},
		department = {数学科学学院},
		major = {概率论与数理统计},
		student-id = {23210180123},
		instructors = {{Élie E. F. Aïdékon\quad 教~授}},%,{某\quad 某\quad 教~授},{某某某\quad 教~授}},
		keywords = {$N$-分支随机游走, 小偏差, $\mathrm L^2$-极限, 渐进速率},
		keywords* = {$N$-BRW, Small Deviations, $\mathrm L^2$-Limit, Asymptotic Speed},
		clc = {O24},
		date = {2024年03月22日}
	}
}

% 序号
\usepackage{enumerate}

% 定义图片文件目录与扩展名
\graphicspath{{figs/}}
\DeclareGraphicsExtensions{.pdf,.eps,.png,.jpg,.jpeg}

% 子图宏包
\usepackage{subcaption}
\usepackage{bicaption}

% 确定浮动对象的位置，可以使用 [H]，强制将浮动对象放到这里（可能效果很差）
\usepackage{float}

% 表格排列
\usepackage{array}

% 固定宽度的表格
\usepackage{tabularx}

% 使用三线表：toprule，midrule，bottomrule。
\usepackage{booktabs}

% 表格中支持跨行
\usepackage{multirow}

% 使用长表格
\usepackage{longtable}

% 附带脚注的表格
\usepackage{threeparttable}

% 附带脚注的长表格
\usepackage{threeparttablex}

% 算法环境宏包
\usepackage[ruled,vlined,linesnumbered]{algorithm2e}
% \usepackage{algorithm, algorithmicx, algpseudocode}

% 代码环境宏包
\usepackage{listings}
\lstdefinestyle{lstStyleCode}{%
  aboveskip         = \medskipamount,
  belowskip         = \medskipamount,
  basicstyle        = \ttfamily\zihao{6},
  commentstyle      = \slshape\color{black!60},
  stringstyle       = \color{green!40!black!100},
  keywordstyle      = \bfseries\color{blue!50!black},
  extendedchars     = false,
  upquote           = true,
  tabsize           = 2,
  showstringspaces  = false,
  xleftmargin       = 1em,
  xrightmargin      = 1em,
  breaklines        = false,
  framexleftmargin  = 1em,
  framexrightmargin = 1em,
  columns           = flexible,
  keepspaces        = true,
  texcl             = true,
  mathescape        = true
}
\lstnewenvironment{codeblock}[1][]{%
  \lstset{style=lstStyleCode,#1}}{}

% 绘图宏包
\usepackage{tikz}
\usetikzlibrary{arrows.meta, shapes.geometric,cd}

% 超链接
\usepackage{xcolor}
\usepackage{hyperref}
\hypersetup{
    colorlinks=true, % 启用彩色链接
    linkcolor=black,  % 普通内部链接的颜色
    citecolor=black,  % 引用链接的颜色
    filecolor=black,  % 打开本地文件的链接颜色
    urlcolor=black    % 链接到URL的颜色
}

% 确保在导言区后期再重新设置一次
\AtEndPreamble{%
    \hypersetup{
        colorlinks=true,     % 启用彩色链接
        linkcolor=red,     % 普通内部链接的颜色
        citecolor=blue,     % 引用链接的颜色
        filecolor=black,     % 本地文件的链接颜色
        urlcolor=black,      % URL 链接的颜色
    }
}

% 页码显示
\usepackage{fancyhdr}
\pagestyle{fancy}
\fancyhf{} 
\fancyfoot[LE,RO]{\thepage} 
\fancypagestyle{plain}{
    \fancyhf{}
    \fancyfoot[LE,RO]{\thepage} 
}
% `L`, `C`, `R`分别对应左、中、右，`E`, `O`对应偶数页和奇数页


% 自定义引用样式，支持可选参数
\NewDocumentCommand{\mycite}{o m}{%
  \IfNoValueTF{#1}
    {\scalebox{1.2}[1.2]{\raisebox{-0.70ex}{\cite{#2}}}}
    { \scalebox{0.9}[0.9]{\raisebox{0.10ex}{\textnormal{[#1, {\scalebox{1.3}[1.3]{\raisebox{-0.80ex}{\citealp{#2}}}}]}}}}}

